\documentclass[a4paper]{article}

\usepackage[margin=2cm]{geometry}
\usepackage{xcolor}
\usepackage{tabularx}
\usepackage{array}
\usepackage[colorlinks=true,linkcolor=blue,urlcolor=blue,citecolor=blue]{hyperref}

\title{Potential Journals for Neurorobotics-Related Papers}
\date{}
\author{Mohammad Rahmani}

\begin{document}
    \maketitle
    
    \renewcommand{\arraystretch}{1.25}
    
    \begin{table}[htbp]
        \centering
        \caption{Potential journals for neurorobotics-related papers}
        \label{tab:neurorobotics_journals}
        
        \begin{tabularx}{\textwidth}{p{0.08\textwidth} p{0.30\textwidth} X}
            \hline
            Priority & Journal & Especially suitable if the paper is about \\
            \hline
            
            1 &
            \href{https://www.frontiersin.org/journals/neurorobotics}{Frontiers in Neurorobotics} &
            Neurorobotics directly: embodied autonomous systems, brain-inspired autonomous systems, neural control, and robot behavior. This is probably the most directly relevant journal. \\
            
            2 &
            \href{https://iopscience.iop.org/journal/1741-2552}{Journal of Neural Engineering} &
            Brain-machine interfaces, neural control, neural prostheses, neurorehabilitation, and neuromorphic engineering. Its scope explicitly includes neurorobotics. \\
            
            3 &
            \href{https://cis.ieee.org/publications/t-cognitive-and-developmental-systems}{IEEE Transactions on Cognitive and Developmental Systems} &
            Cognitive robotics, developmental robotics, autonomous and evolutionary robotics, and computational neuroscience. It is suitable when the paper focuses on learning, cognition, or embodiment in robots. \\
            
            4 &
            \href{https://ieeexplore.ieee.org/xpl/RecentIssue.jsp?punumber=7333}{IEEE Transactions on Neural Systems and Rehabilitation Engineering} &
            Brain-computer interfaces, EMG/EEG, motor control, assistive devices, rehabilitation robotics, and neuroprosthetics. It is suitable when the paper has a medical or neurotechnological focus. \\
            
            5 &
            \href{https://link.springer.com/journal/12984}{Journal of NeuroEngineering and Rehabilitation} &
            Exoskeletons, wearable sensing, neurorehabilitation, and robot-assisted rehabilitation. It is suitable when the paper concerns therapy, stroke rehabilitation, or assistive robotics. \\
            
            6 &
            \href{https://www.ieee-ras.org/publications/ra-l/}{IEEE Robotics and Automation Letters} &
            General robotics, including control, learning, perception, and automation. It is suitable when the robotics contribution is stronger than the neuroscience contribution. \\
            
            7 &
            \href{https://www.science.org/journal/scirobotics}{Science Robotics} &
            High-impact general robotics papers. It is suitable only if the contribution is clearly novel and experimentally strong. \\
            
            8 &
            \href{https://iopscience.iop.org/journal/2634-4386}{Neuromorphic Computing and Engineering} &
            Spiking neural networks, neuromorphic hardware, brain-inspired computing, and neural controllers for robots. \\
            
            9 &
            \href{https://link.springer.com/journal/422}{Biological Cybernetics} &
            Theoretical models of motor control, sensorimotor systems, information processing, and biological control. It is suitable when the paper is more mathematical or theoretical. \\
            
            10 &
            \href{https://iopscience.iop.org/journal/1748-3190}{Bioinspiration \& Biomimetics} &
            Bio-inspired robotics, biomimetic control, and robot design inspired by biological systems. \\
            
            11 &
            \href{https://www.frontiersin.org/journals/computational-neuroscience}{Frontiers in Computational Neuroscience} &
            Computational neuroscience and theoretical modeling of brain function. It is suitable when the robot mainly serves as a simulation platform or model. \\
            
            12 &
            \href{https://www.sciencedirect.com/journal/neural-networks}{Neural Networks} &
            Neural models, artificial neural networks, and learning systems. It is suitable when the main contribution is algorithmic or neural-network-oriented. \\
            
            \hline
        \end{tabularx}
    \end{table}

\end{document}